\documentclass[12pt,a4paper,oneside]{article}
\usepackage[brazil]{babel}
\usepackage[utf8]{inputenc}
\usepackage{olymp}

\setcounter{problem}{5}

\begin{document}

\begin{problem}{Variante}{variante}{2}

Preocupados com novas variantes do novo coronavírus, mais contagiosas, um grupo de pesquisadores está fazendo estudos para simular o crescimento do número de infectados quando não são tomadas medidas de prevenção, como distanciamento social e isolamento de pessoas infectadas.

Suponha, por exemplo, que cada pessoa infectada seja capaz de infectar outras 2 por semana. Para simplificar, não vamos considerar recuperação, ou seja, cada infectado continua infectando novas pessoas nessa mesma taxa indefinidamente. Assim, partindo de apenas 1 pessoa infectada, em 1 semana serão 3 pessoas (ela e mais duas); em 2 semanas já serão 9; em 3 semanas 27, e assim por diante. É possível que rapidamente a população infectada chegue a 1.000.000.000 de pessoas!

Faça um programa para, sabendo o tamanho da população inicial de pessoas infectadas por uma nova variante e a transmissibilidade da variante, que indica quantas pessoas cada uma consegue infectar por semana, calcule em quanto tempo a população ultrapassa 1.000.000.000 de infectados.

%\Notes
%
%Observações, se tiver.

\InputFile

A entrada é composta por dois inteiros, $P_0$ e $G$, que representam respectivamente a população inicial de pessoas infectadas (1, no exemplo acima) e a taxa de transmissibilidade, ou seja, quantas novas pessoas são infectadas por pessoa infectada a cada semana (2, no exemplo acima). Restrições: $1 \leq P_0, G \leq 1000$.


\OutputFile

Seu programa deve gerar apenas uma linha de saída, contendo um único inteiro que indica quantas semanas serão necessárias para a população atingir no mínimo 1.000.000.000 de infectados.
% \Notes
% Preste atenção aos tipos das variáveis, pois $20! \approx 2^{61}$.

\Example

\begin{example}
\exmp{
1 2
}{
19
}%
\end{example}

\begin{example}
\exmp{
1 9
}{
9
}%
\end{example}

\begin{example}
\exmp{
3 9
}{
9
}%
\end{example}

\begin{example}
\exmp{
6 13
}{
8
}%
\end{example}

%Comentários sobre o exemplo, se necessário.

\end{problem}

\end{document}
